
\textbf{Burnside's Lemma/Polya Enumeration}
\begin{lstlisting}
using Permutation = vector<int>;
void operator*=(Permutation& p, Permutation
const& q) {
  Permutation copy = p;
  for (int i = 0; i < p.size(); i++)
    p[i] = copy[q[i]];
}
int count_cycles(Permutation p) {
  int cnt = 0;
  for (int i = 0; i < p.size(); i++) {
    if (p[i] != -1) {
      cnt++;
      for (int j = i; p[j] != -1;) {
        int next = p[j];
        p[j] = -1;
        j = next;
      }
    }
  }
  return cnt;
}
int solve(int n, int m) {
  Permutation p(n*m), p1(n*m), p2(n*m),
  p3(n*m);
  for (int i = 0; i < n*m; i++) {
    p[i] = i;
    p1[i] = (i % n + 1) % n + i / n * n;
    p2[i] = (i / n + 1) % m * n + i % n;
    p3[i] = (m - 1 - i / n) * n + (n - 1 - i % n);
  }
  set<Permutation> s;
  for (int i1 = 0; i1 < n; i1++) {
    for (int i2 = 0; i2 < m; i2++) {
      for (int i3 = 0; i3 < 2; i3++) {
        s.insert(p);
        p *= p3;
      }
      p *= p2;
    }
    p *= p1;
  }
  int sum = 0;
  for (Permutation const& p : s) {
    sum += 1 << count_cycles(p);
  }
  return sum / s.size();
}
\end{lstlisting}

\textbf{Exclusion DP}
There is only difference to calculate $f[i]$ for different problems,

for ordered pair means $(a, b)$ and $(b, a)$ are not same then $f[i] = d[i] \cdot d[i]$;

for unordered pair means $(a, b)$ and $(b, a)$ are same then

-for repeative $\to f[i] = (d[i]+1)C2 \to d[i] \cdot (d[i]+1)/2$;

-for non-repeative $\to f[i] = d[i]C2 \to d[i] \cdot (d[i]-1)/2$;

for $k$-tuples:

-for ordered tuples : $f[i] = (d[i])^k$;

-for unordered tuples :

--for non-repeative $\to f[i] = d[i]Ck$

--for repeative $\to f[i] = (d[i]+k-1)Ck$

here $d[i]$ = number of elements divisible by $i$.

and $C$ denotes combination

\textbf{Bishop Placement}
/* Number of ways to place $K$ bishops on an $N \times N$ chessboard so that no two bishops attack each other. */
\begin{lstlisting}
int bishop_placements(int N, int K)
{
  if (K > 2 * N - 1)
    return 0;
  vector<vector<int>> D(N * 2, vector<int>(K +
  1));
  for (int i = 0; i < N * 2; ++i)
    D[i][0] = 1;
  D[1][1] = 1;
  for (int i = 2; i < N * 2; ++i)
    for (int j = 1; j <= K; ++j)
      D[i][j] = D[i-2][j] + D[i-2][j-1] *
      (squares(i) - j + 1);
  int ans = 0;
  for (int i = 0; i <= K; ++i)
    ans += D[N*2-1][i] * D[N*2-2][K-i];
  return ans;
}
\end{lstlisting}

\textbf{Handling Precision (EPS)}
\begin{lstlisting}
const double EPS = 1e-9;
bool equal(double a, double b) { return abs(a - b)
< EPS; }
bool lessThan(double a, double b) { return a < b - EPS; }
bool greaterThan(double a, double b) { return a >
b + EPS; }
\end{lstlisting}

\textbf{Point and Vector}
\begin{lstlisting}
const double EPS = 1e-9;
struct Point {
  double x, y;
  Point operator+(const Point& other) const {
    return {x + other.x, y + other.y}; }
  Point operator-(const Point& other) const {
    return {x - other.x, y - other.y}; }
  Point operator*(double s) const { return {x * s,
  y * s}; }
  Point operator/(double s) const { return {x / s,
  y / s}; }
  bool operator<(const Point& other) const {
    if (abs(x - other.x) > EPS) return x < other.x;
    return y < other.y - EPS;
  }
  bool operator==(const Point& other) const {
    return abs(x - other.x) < EPS && abs(y -
    other.y) < EPS;
  }
};
typedef Point Vector;
double dot(Vector a, Vector b) { return a.x * b.x +
a.y * b.y; }
double norm_sq(Vector a) { return dot(a, a); }
double length(Vector a) { return
sqrt(norm_sq(a)); }
double cross(Vector a, Vector b) { return a.x * b.y
- a.y * b.x; }
\end{lstlisting}

\begin{lstlisting}
// Returns: > 0 for Left Turn, < 0 for Right Turn,
// 0 for Collinear
double orientation(Point p, Point a, Point b) {
  return cross(a - p, b - p);
}
double distToLine(Point p, Point a, Point b) {
  return abs(cross(p - a, b - a)) / length(b - a);
}
double distToSegment(Point p, Point a, Point b) {
  if (dot(p - a, b - a) < 0) return length(p - a);
  if (dot(p - b, a - b) < 0) return length(p - b);
  return distToLine(p, a, b);
}
double polygonArea(const vector<Point>& p) {
  double area = 0;
  for (int i = 0; i < (int)p.size(); i++) {
    area += cross(p[i], p[(i + 1) % p.size()]);
  }
  return abs(area) / 2.0;
}
\end{lstlisting}

\begin{lstlisting}
void convex_hull(vector<pt>& a, bool
include_collinear = false) {
  pt p0 = *min_element(a.begin(), a.end(), [](pt
  a, pt b) {
    return make_pair(a.y, a.x) < make_pair(b.y,
    b.x);
  });
  sort(a.begin(), a.end(), [&p0](const pt& a, const
  pt& b) {
    int o = orientation(p0, a, b);
    if (o == 0)
      return (p0.x-a.x)*(p0.x-a.x) +
      (p0.y-a.y)*(p0.y-a.y)
      < (p0.x-b.x)*(p0.x-b.x) +
      (p0.y-b.y)*(p0.y-b.y);
    return o < 0;
  });
  if (include_collinear) {
    int i = (int)a.size()-1;
    while (i >= 0 && collinear(p0, a[i], a.back()))
      i--;
    reverse(a.begin()+i+1, a.end());
  }
  vector<pt> st;
  for (int i = 0; i < (int)a.size(); i++) {
    while (st.size() > 1 && !cw(st[st.size()-2],
    st.back(), a[i], include_collinear))
      st.pop_back();
    st.push_back(a[i]);
  }
  if (include_collinear == false && st.size() == 2
  && st[0] == st[1])
    st.pop_back();
  a = st;
}
